\documentclass[preview]{standalone}

\usepackage{amsmath}
\usepackage{amssymb}
\usepackage{stellar}
\usepackage{definitions}

\begin{document}

\id{finite-abelian-groups-exercises}
\genpage

\section{Exercises}

\begin{snippetexercise}{finite-abelian-groups-cyclic-products-exercise}{Cyclic groups and direct products}
    Determine which of the following groups are isomorphic:
    \[
        C_4\cartesianprod C_9,
        \quad
        C_6\cartesianprod C_6,
        \quad
        C_8\cartesianprod C_3,
        \quad
        C_9\cartesianprod C_4,
        \quad
        C_6\cartesianprod C_4,
        \quad
        C_{24}.
    \]
\end{snippetexercise}

\begin{snippetsolution}{finite-abelian-groups-cyclic-products-solution}{Cyclic groups and direct products}
    The product \(C_m\cartesianprod C_n\) is cyclic of order \(mn\)
    \ifandonlyif \(\gcd(m,n)=1\). Hence
    \[
        C_4\cartesianprod C_9
        \groupisomorphic C_{36}
        \groupisomorphic C_9\cartesianprod C_4
    \]
    and
    \[
        C_8\cartesianprod C_3\groupisomorphic C_{24}.
    \]

    The group \(C_6\cartesianprod C_6\) is not cyclic: all its elements
    have order at most \(6\), while the group has order \(36\). Finally,
    \[
        C_6\cartesianprod C_4
        \groupisomorphic
        C_3\cartesianprod C_2\cartesianprod C_4
        \groupisomorphic
        C_2\cartesianprod C_{12}.
    \]
    It has order \(24\), but its elements have order at most \(12\), so
    it is not isomorphic to \(C_{24}\).

    Therefore the isomorphism classes in the list are
    \[
        \{C_4\cartesianprod C_9,C_9\cartesianprod C_4\},
        \qquad
        \{C_8\cartesianprod C_3,C_{24}\},
    \]
    \[
        \{C_6\cartesianprod C_6\},
        \qquad
        \{C_6\cartesianprod C_4\}.
    \]
\end{snippetsolution}

\begin{snippetexercise}{finite-abelian-groups-orders-divide-twenty-eight-exercise}{Groups whose element orders divide 28}
    Determine all finite \abeliangroup[abelian groups] \(G\) such that
    \[
        30\leq |G|\leq60
    \]
    and
    \[
        \elementperiod{g}\divides28
        \qquad
        \text{for every }g\in G.
    \]
\end{snippetexercise}

\begin{snippetsolution}{finite-abelian-groups-orders-divide-twenty-eight-solution}{Groups whose element orders divide 28}
    If a prime \(p\) divides \(|G|\),
    \snippetref[cauchy-lemma][Cauchy's theorem] gives an element of
    order \(p\). Hence \(p\divides28\), so only \(2\) and \(7\) may
    divide \(|G|\). The orders in the prescribed interval of the form
    \(2^a7^b\) are
    \[
        32,
        \qquad
        49,
        \qquad
        56.
    \]

    If \(|G|=32=2^5\), the exponent of \(G\) must divide \(4\). The
    partitions of \(5\) having no part larger than \(2\) give
    \[
        C_4\cartesianprod C_4\cartesianprod C_2,
        \qquad
        C_4\cartesianprod C_2\cartesianprod C_2\cartesianprod C_2,
        \qquad
        C_2^5.
    \]

    If \(|G|=49\), the two possibilities are \(C_{49}\) and
    \(C_7\cartesianprod C_7\). The first contains elements of order
    \(49\nmid28\), so only
    \[
        C_7\cartesianprod C_7
    \]
    remains.

    If \(|G|=56=2^3\cdot7\), the \(7\)-primary component is \(C_7\).
    The \(2\)-primary component cannot be \(C_8\), while the other two
    possibilities yield
    \[
        C_4\cartesianprod C_2\cartesianprod C_7
        \groupisomorphic C_2\cartesianprod C_{28}
    \]
    and
    \[
        C_2\cartesianprod C_2\cartesianprod C_2\cartesianprod C_7
        \groupisomorphic C_2\cartesianprod C_2\cartesianprod C_{14}.
    \]
    These six displayed groups are all the solutions.
\end{snippetsolution}

\end{document}
